% Suivi et navigation embarqués
\section{Suivi et navigation embarqués}
\label{sec:navigation}

\subsection{Suivi mono-objet sur Jetson}

Le groupe Vision4Robotics (Tongji) publie trois trackers \glslink{siamtracker}{siamois} à \gl{backbone}{backbone} \glslink{alexnet}{AlexNet},
testés à bord d'un drone sur \glslink{agx-xavier}{Jetson AGX Xavier}, tous sans TensorRT et sans mesure de
puissance.

\begin{table}[!htbp]
\centering\footnotesize
\caption{Trackers aériens mesurés à bord (gamme Jetson). Aucune puissance ni mode de
puissance n'est rapporté.}
\label{tab:trackers}
\begin{tabularx}{\textwidth}{@{}P{3.4cm}LrP{2.8cm}@{}}
\toprule
Tracker & Précision sur benchmark (bureau) & \glslink{fps}{FPS} à bord & \glslink{fps}{FPS} bureau \\
\midrule
\glslink{siamapnpp}{SiamAPN++} \citep{arxiv-2106.08816-siamapn-siamese-attentional-aggregation-network}\pp &
  UAV20L : précision 0,736, \glslink{auc}{AUC} 0,560 \src{Fig.~4a}\med & 34,9 \src{§V}\med & \nr \\
\glslink{hift}{HiFT} \citep{arxiv-2108.00202-hift-hierarchical-feature-transformer-aerial}\pp &
  UAV20L : précision 0,763 \src{Table~2} & 31,2 \src{§5}\med & 127 (\gl{GPU}{GPU} non précisé) \\
\glslink{tctrack}{TCTrack} \citep{arxiv-2203.01885-tctrack-temporal-contexts-aerial-tracking}\pp &
  UAVTrack112\_L : succès 0,582 \src{§4.2}\med & > 27 \src{§5}\med & 125,6 (« PC ») \\
\bottomrule
\end{tabularx}
\end{table}

Le passage du bureau à la \glslink{agx-xavier}{Xavier} divise le débit par quatre environ, sans analyse du
goulot \deduit. \glslink{tctrack}{TCTrack} apporte une observation utile au choix d'un backbone : sur \glslink{agx-xavier}{Xavier},
\glslink{alexnet}{AlexNet} prend 3,4\,\glslink{u-ms}{ms} contre 13,7\,ms pour \glslink{mobilenetv2}{MobileNetV2}, attribués au coût d'accès mémoire
des réseaux « mobiles » \citep[Table~1]{arxiv-2203.01885-tctrack-temporal-contexts-aerial-tracking}.
Aucun tracker siamois ou à \glslink{attention}{Transformer} n'est porté sur \gl{FPGA}{FPGA} ou \gl{MCU}{MCU} dans le corpus, et
aucun n'étudie la \gl{quantif}{quantification}. Du côté multi-objets, les trackers par association de
positions sont quasi gratuits : SORT tourne à 245\,\glslink{fps}{FPS} sur \gl{CPU}{CPU} dans UAVDT
\citep[Table~3]{arxiv-1804.00518-unmanned-aerial-vehicle-benchmark-object}\med. Le coût
embarqué du suivi est donc dans le détecteur \deduit.

\subsection{Vol agile : la référence Jetson de Zurich}

Deux travaux de l'université de Zurich fixent les ordres de grandeur de la \glslink{gamme}{gamme Jetson}
pour un drone agile de moins d'un kilogramme.
\begin{itemize}
  \item \citet{arxiv-2110.05113-learning-high-speed-flight-wild} (Science Robotics)
    apprennent une politique de \gl{bout}{bout en bout} à partir d'images de profondeur. Sur \glslink{tx2}{Jetson
    TX2}, la latence est de 41,6\,\glslink{u-ms}{ms}, dont 38,9 d'inférence \src{Table~1}. En forêt, le
    drone réussit 100\,\% des vols à 5\,m/s et 60\,\% à 10\,m/s \src{§2-A}. Le papier
    relie explicitement latence et \glslink{vsure}{vitesse sûre} maximale \src{§S5}.
  \item \citet{2023-kaufmann-champion-level-drone-racing-deep} (Nature) battent des
    champions humains de course FPV avec des capteurs et un calcul embarqués. Le détecteur
    de portes (un petit U-Net) prend 40\,ms en TensorRT \gl{FP32}{FP16} sur le GPU du \glslink{tx2}{TX2}, la
    politique 8\,ms sur son CPU (Methods)\med. La vision est le goulot, et la caméra à
    30\,\glslink{u-hz}{Hz} est un facteur limitant face aux 120\,Hz des pilotes (Discussion).
\end{itemize}
Aucun des deux ne mesure la puissance du calculateur. Le survey de
\citet{arxiv-2301.01755-autonomous-drone-racing-survey} affirme que la famille \glslink{jetson}{Jetson} est
le plus petit matériel suffisant pour des réseaux complexes, ce qui fixe des châssis de
6 pouces \src{§VII-C} : il ignore les nano-drones et ne discute pas le FPGA. Il cite comme
question ouverte la détection d'un adversaire depuis la caméra embarquée (\emph{multiplayer
racing}, §VIII-C), c'est-à-dire un cas \gl{airair}{air $\rightarrow$ air}.

\paragraph{Politiques apprises en simulation.} Deux \gl{preprint}{preprints} récents poussent la
politique pixels $\rightarrow$ commandes (poussée collective et vitesses angulaires) plus
loin, mais sans aucune mesure de calcul à bord. \citet{arxiv-2501.14377-dream-fly-model-reinforcement-learning}\pp{}
entraînent une politique de course par \glslink{dreamerv3}{DreamerV3} (apprentissage par renforcement avec modèle
du monde) sur des images RGB 64$\times$64 \src{§IV} ; PPO et SAC n'apprennent aucun vol utile
dans les mêmes conditions (Fig.~3, sans valeur numérique). L'entraînement prend environ
240\,h sur un \gl{GPU}{GPU} de bureau Quadro RTX~8000 \src{§IV-A}, et le vol « réel », jusqu'à 9\,m/s
(résumé), est un \gl{hil}{hardware-in-the-loop} : images rendues par simulateur, commandes
calculées sur un ordinateur déporté \src{§V-A}. Latence, débit, puissance et taille du
réseau embarqué : \nr. \citet{arxiv-2604.12916-e2e-fly-integrated-training-deployment}\pp{}
proposent une chaîne complète d'entraînement (simulation différentiable, PPO ou rétropropagation
dans le temps), de validation et de déploiement, avec un quadrirotor embarquant un \glslink{n100}{Intel N100}
(x86 basse consommation, \glslink{gamme}{gamme} non mesurée) \src{§VI-A}. Le seul chiffre de calcul est une
latence de communication d'environ 90\,\glslink{u-ms}{ms} en commande déportée, ramenée sous 30\,ms avec le
calcul à bord \src{§VI-C2}. Les 100\,\% de réussite en course avec obstacles sont obtenus en
simulation \src{§VII-E} ; le transfert réel n'est montré qu'en figures, sans taux de succès
ni baseline publiée. Ces deux travaux confirment que la politique de contrôle elle-même
n'est pas le goulot \deduit, mais ni l'un ni l'autre ne donne la latence d'inférence de sa
politique sur le calculateur de bord.

\subsection{Suivi de cible et navigation sémantique sur Jetson Orin}

Quatre \gl{preprint}{preprints} de 2025--2026 font tourner une chaîne complète perception
$\rightarrow$ estimation $\rightarrow$ contrôle sur un module \glslink{orin}{Jetson Orin}, en vol réel
(tableau~\ref{tab:orin-nav}). Aucun ne rapporte de puissance ni de mode de puissance :
ils se situent dans la \glslink{gamme}{gamme Jetson} (5 à 30\,\glslink{u-w}{W} selon le module et le mode) sans plus de
précision \deduit. Le seul repère de puissance \glslink{orin}{Orin} du corpus porte sur une autre charge,
\glslink{yolov8n}{YOLOv8n} \gl{INT8}{INT8} (10,3\,W sur \glslink{orin-nx}{Orin NX}, 7,5\,W sur \glslink{orin-nano}{Orin Nano}, par
\gl{telemetrie}{télémétrie}) \citep[Table~4]{arxiv-2502.15737-performance-analysis-you-only-look} :
un ordre de grandeur, pas une mesure de ces systèmes.

\begin{table}[!htbp]
\centering\footnotesize
\caption{Chaînes de navigation et de suivi mesurées à bord sur Jetson Orin (tous
preprints). Puissance \texttt{n/r} partout ; résultats de mission sur les essais propres
aux auteurs, non comparables entre eux.}
\label{tab:orin-nav}
\begin{tabularx}{\textwidth}{@{}P{2.6cm}P{2.2cm}LL@{}}
\toprule
Système & Calculateur & Calcul mesuré à bord & Résultat de mission \\
\midrule
\glslink{yopov2}{YOPOv2-Tracker} \citep{arxiv-2505.06923-yopov2-tracker-end-end-agile}\pp &
  \glslink{orin-nx}{Orin NX} & perception $\rightarrow$ action 8,2\,\glslink{u-ms}{ms} en moyenne \src{§IV-B-3}\med ; précision
  numérique \nr & suivi d'un humain à 6\,m/s en forêt, 5,6\,m/s entre bâtiments
  \src{§IV-B-3} \\
NOVA \citep{arxiv-2506.18689-nova-navigation-object-centric-visual}\pp &
  \glslink{orin-nx}{Orin NX} 16\,Go & \glslink{yolov11s}{YOLOv11-s} \gl{FP32}{FP16} TensorRT : 92\,\glslink{u-hz}{Hz} \src{§IV-B}\med ; complétion de
  profondeur (\glslink{depthanythingv2}{DepthAnythingV2} ViT-S) : 31,6\,\glslink{u-ms}{ms} \src{§IV-C}\med & détections 93,1 à
  98,6\,\% selon le scénario, cible à plus de 53\,km/h en forêt \src{§V-E}\med \\
NGPS \citep{arxiv-2607.18936-ngps-gps-denied-aerial-geo}\pp &
  \glslink{orin-nx}{Orin NX} & appariement satellite 2,6\,\glslink{u-hz}{Hz} sur GPU, 0,5\,Hz sur \gl{CPU}{CPU} \src{§IV-D}\med &
  \glslink{rmse}{RMSE} de position 2,94\,m contre 10,37\,m en VIO seul \src{§IV-B, §IV-E4}\med \\
\glslink{singer}{SINGER} \citep{arxiv-2509.18610-singer-onboard-generalist-vision-language}\pp &
  \glslink{orin-nano}{Orin Nano} 8\,Go & \glslink{clipseg}{CLIPSeg} (ViT-B/16, 86\,M param.) 12\,\glslink{u-hz}{Hz} en ONNX Runtime, 3\,Hz sans
  optimisation \src{§V}\med & 76,67\,\% des essais à moins de 1\,m du but (30 essais),
  contre 53,33\,\% pour la baseline \src{§VI-D}\med \\
\bottomrule
\end{tabularx}
\end{table}

\paragraph{Suivi agile de bout en bout.} \citet{arxiv-2505.06923-yopov2-tracker-end-end-agile}\pp{}
remplacent la chaîne détection $\rightarrow$ carte $\rightarrow$ planification par un seul
réseau convolutionnel (\glslink{resnet}{ResNet-18}, entrée RGB-D 160$\times$96) qui prédit directement des
primitives de trajectoire, à la manière des boîtes d'ancrage d'un détecteur \src{§III-C}.
Sur un quadrotor de 850\,g, la latence perception $\rightarrow$ action est de 8,2\,\glslink{u-ms}{ms} sur
\glslink{orin-nx}{Orin NX} \src{§IV-B-3}, soit près de cinq fois moins que l'inférence de
\citet{arxiv-2110.05113-learning-high-speed-flight-wild} sur \glslink{tx2}{TX2} \deduit{} (calculateurs et
réseaux différents). Les 3\,ms annoncés en simulation sont mesurés sur un GPU de bureau \glslink{rtx30mid}{RTX~3060}
\src{§IV-B-1}. Les baselines (Elas.~Tracker, Vis.~Tracker) sont conçues pour
2 à 2,5\,m/s et les taux de succès ne sont donnés qu'en figure, sur 10 essais par vitesse :
comparaison favorable à la méthode \deduit.

\paragraph{Suivi à haute vitesse sans GPS.} \citet{arxiv-2506.18689-nova-navigation-object-centric-visual}\pp{}
suivent une personne ou un véhicule sans capture de mouvement ni localisation globale, dans
le repère de la cible : détecteur \glslink{yolov11s}{YOLOv11-s} avec zoom adaptatif, complétion de profondeur,
\glslink{ukf}{filtre de Kalman \emph{unscented}} et commande prédictive non linéaire, le tout sur \glslink{orin-nx}{Orin NX}
\src{§III}. Le chiffre utile pour le calcul est l'écart entre modules : la complétion de
profondeur (31,6\,\glslink{u-ms}{ms}) est environ trois fois plus lente que le détecteur
(92\,\glslink{u-hz}{Hz}, soit 10,9\,ms) \deduit. Sur ce pipeline, le coût n'est donc pas dans le détecteur,
contrairement au suivi multi-objets ci-dessus. Le « \glslink{detrate}{taux de détection} » n'est pas défini
contre une vérité terrain : il ne se compare ni à une \glslink{map}{mAP} ni à un \glslink{iou}{IoU}. Un seul vol par
scénario, aucune baseline, aucune latence de \gl{bout}{bout en bout}.

\paragraph{Navigation guidée par le langage.} \citet{arxiv-2509.18610-singer-onboard-generalist-vision-language}\pp{}
pilotent un drone de 5 pouces vers un objet désigné en langage naturel, avec tout le calcul
sur \glslink{orin-nano}{Jetson Orin Nano}. La perception est un modèle de fondation, \glslink{clipseg}{CLIPSeg} (segmentation
conditionnée par le texte), dont la carte de chaleur alimente une petite politique
visuomotrice à 20\,\glslink{u-hz}{Hz} \src{§V}. \glslink{clipseg}{CLIPSeg} est le goulot : 12\,Hz au mieux après conversion
ONNX, sous la fréquence de la commande \src{§V, §VI-C}. Sur 30 essais en environnement et
requêtes inconnus, 76,67\,\% atteignent le but à moins de 1\,m (66,67\,\% sans cap de
référence externe) \src{§VI-D}\med ; la baseline est une implantation minimale écrite par
les auteurs. Pas d'évitement d'obstacles, ni de puissance, ni de latence mesurée. C'est le
seul modèle de type vision--langage exécuté à bord dans cette section.

\paragraph{Ce que ces quatre travaux changent.} Ils déplacent la référence de la \glslink{gamme}{gamme
Jetson} du \glslink{tx2}{TX2} vers l'\glslink{orin}{Orin}, et montrent qu'en suivi et en navigation sémantique le goulot
est un réseau de perception dense (profondeur, segmentation, appariement), pas le
détecteur ni la politique \deduit. Ils reproduisent aussi les manques de la littérature
Zurich : aucune puissance, aucune autonomie, aucun lien chiffré entre cadence de
perception et \glslink{vsure}{vitesse sûre}, aucune \gl{quantif}{quantification} au-delà du FP16, aucun portage
\gl{FPGA}{FPGA} ou \gl{MCU}{MCU}.

\subsection{Localisation et plateformes hétérogènes}

\citet{2021-guo-fpga-uav-heterogeneous-computing-platform} proposent de réunir contrôle
de vol (\gl{PID}{PID}) et prétraitement sur un \glslink{xc7z020}{Zynq-7020}, avec une \glslink{agx-xavier}{Jetson AGX Xavier} pour
l'inférence. C'est une proposition sans aucune mesure de latence, de masse ou de puissance,
à citer avec prudence \src{§III--IV}. \citet{arxiv-2307.10255-land-localize-infrastructure-free-scalable}\pp{}
localisent un essaim de nano-drones par \gl{UWB}{UWB}, sans vision ; le module radio seul consomme
environ 350\,\glslink{u-mw}{mW} par drone \src{§V-C}, un ordre de grandeur à garder face au budget de
calcul d'un nano-drone.

La localisation visuelle absolue sans GNSS n'a qu'un représentant.
\citet{arxiv-2607.18936-ngps-gps-denied-aerial-geo}\pp{} (auteur unique ; accepté à IROS 2026
selon le PDF, non vérifié) apparient l'image d'une caméra orientée vers le bas avec une
image satellite (\glslink{superpoint}{SuperPoint} + \glslink{lightglue}{LightGlue}) et fusionnent le résultat avec un VIO et l'IMU,
à 60--150\,m d'altitude \src{§III}. Sur \glslink{orin-nx}{Orin NX}, l'appariement tourne à 2,6\,\glslink{u-hz}{Hz} sur GPU et
0,5\,Hz sur CPU, \glslink{lightglue}{LightGlue} prenant 82\,\% du temps GPU \src{§IV-D}\med. L'erreur de position
fusionnée est de 2,94\,m contre 10,37\,m pour le VIO seul \src{§IV-B, §IV-E4}\med, sur cinq
vols à 2--4\,m/s ; l'\glslink{ate}{ATE} atteint 6,04\,m au pire à 150\,m d'altitude \src{§IV-E2}. La comparaison à
un autre VIO vient d'un autre article et d'autres données (non appariée), et puissance,
mémoire et latence de \gl{bout}{bout en bout} sont \nr. L'appariement par \glslink{attention}{attention} de
\glslink{lightglue}{LightGlue}, goulot mesuré, n'a aucun équivalent \gl{FPGA}{FPGA} ou \gl{NPU}{NPU} dans le corpus.
